\section{延伸：纽结同调的范畴化}
\label{sec:categorification}

这一节超出本文范围，但用户点名要作为延伸。问题：Witten 1989 与 Khovanov 同调是什么关系？

一句话答案：\textbf{不是推论关系，但 1989 年的 CS 是“物理来源”的起点}。Khovanov 同调把 Jones 多项式从“一个数”升级成“一族双分次向量空间”，Jones 是它的分次 Euler 特征；而 Witten 在 2011 年把这件事的物理来源放到了高一维（4d/5d），并且明说“范畴化”在他的口径里就是“能从高一维理论导出”。

\subsection{范畴化：数 $\to$ 复形}

Jones 多项式是一个数（Laurent 多项式）。Khovanov 同调把每个链环 $L$ 换成一个双分次同调群 $H^{i,j}(L)$，使得（在适当归一化下；不同作者的 $q$ 与 $q+q^{-1}$ 因子约定不同）
\begin{equation}
\sum_{i,j}(-1)^i q^j\dim H^{i,j}(L)=J_L(q).
\end{equation}
即 Jones 多项式是它的 \textbf{Euler 特征} \ours。所以 Khovanov \textbf{没有证明} Witten 的 (1.6)：他假定 Jones 多项式已知，构造一个以它为 Euler 特征的同调。这与本文不在同一层。

升级换来的东西有两样，都是 1989 年没有的：
\begin{itemize}[leftmargin=1.6em]
\item \textbf{函子性}：链环协边 $\Sigma:K\to K'$ 给出线性映射 $\Phi_\Sigma:\mathrm{Kh}(K)\to\mathrm{Kh}(K')$，拼接相乘；1989 的 TQFT 只有“数 $\to$ 数”；
\item \textbf{解释整数性}：Jones 的系数为什么是整数？CS 路径积分给 Laurent 多项式，但不说这件事；Khovanov 把系数解释成分次向量空间的维数——Witten 2011 把这句写成 §1.1.2 的标题 “What Chern--Simons Theory Doesn't Explain” \src{Witten 2011 p.8}。
\end{itemize}

\subsection{物理来源是后世的，且“范畴化 = 升一维”}

本文给的是 Jones 的 3d CS 描述。Khovanov 同调的物理来源出现在 2007--2011 年：
\begin{itemize}[leftmargin=1.6em]
\item Witten 2011：D3 膜终止于带 $\theta$ 角的 NS5 膜，两支 S/T 对偶给出 Jones 与 Khovanov 的流形无关描述（4/5 维椭圆方程解的计数），arXiv:1101.3216；
\item Kapustin--Witten 2007：4d $\cN=4$ SYM 的拓扑扭曲族，提供 Wilson--'t Hooft 与面算子框架；它讲几何 Langlands，Khovanov 不是它的推论，只是同一框架的另一应用，arXiv:hep-th/0604151。
\end{itemize}

Witten 自己的点题句值得原样记住：
\begin{quote}
``the fact that the four-dimensional equations (5.38) can be `lifted' \dots\ to five dimensions is part of the reason that Chern--Simons gauge theory can be `categorified,' which is just a fancy way to say that it can be derived from a theory in one dimension higher.''\\[2pt]
\hfill --- Witten 2011, p.107
\end{quote}

这是\textbf{并列的两套物理实现}（3d CS 对 Jones；4d/5d 对 Khovanov），不是同一套的延伸。任何“Witten 1989 已经预示 Khovanov”的说法都是时代错置。

\subsection{本文有没有范畴化暗示：没有}

全文没有 homology、chain complex、categorification 一类构造，不变量始终是数 $Z(M;\text{links})$。唯一擦边处是 p.372：解带标记点约束需要“$q$ 数联络”，其和乐 “would presumably be an element of a `quantum group', not an ordinary classical group. Indeed, it seems likely that\dots” \src{372}——“presumably / seems likely” 双重 hedging，且说的是量子群，不是同调。

但本文有“半截”结构：对曲面 $\Sigma$ 给向量空间 $\cH_\Sigma$（§3），对三维流形给数——已是 Atiyah TQFT 的雏形，只是取值为向量空间而非复形，缺 $\ZZ$-分次与微分。范畴化要做的正是把 $\cH_\Sigma$ 换成复形、把 $Z(M)$ 换成上同调。

\subsection{要把本文理论范畴化，缺什么}

五条开放问题 \ours：
\begin{enumerate}[leftmargin=2em]
\item \textbf{一维提升}：需要 4d（或 5d/6d）理论，其边界约化给 CS；3d 理论自身没有“多出来的一维”放同调方向。
\item \textbf{双分次}：需要 $Q^2=0$ 与 $U(1)\times U(1)$ 给 $(i,j)$ 分次；纯 CS 只有规范约束，没有 $R$-对称性。
\item \textbf{标架/半形式反常}：CS 配分函数依赖标架（更一般地 two-framing），换标架乘 $q$ 的幂；1989 用 $\eta$ 不变量 (2.16)(2.17) 处理相位。要得到良定义的 $\ZZ$-分次复形，必须把反常典范化。
\item \textbf{表示范畴升级}：Wilson 线只带一个可积表示，对应一个模张量范畴；范畴化要求把它换成 2-范畴。一般 $G$ 未知，$A$ 型有 Khovanov--Rozansky/Soergel 路线。
\item \textbf{level 与 $q$ 的错位}：1989 的不变量定义在 $q=e^{2\pi\ii/(k+N)}$（固定 level），Khovanov 是 $q$ 的多项式（所有 level 同时）；“逐 level 范畴化”与 Khovanov 同调不是同一件事。
\end{enumerate}

T6 的真空也在这里：1989 的“CS $\to$ Jones”与辛几何路线的“$\mathrm{Kh}_{\rm symp}\to$ Khovanov”之间，目前没有一条已知定理把 CS 的 Wilson 线直接变成 Khovanov 的链复形。本地线索：\texttt{functoriality\_work/work/通读笔记.md:55}（在 \texttt{C:\textbackslash Users\textbackslash admin\textbackslash Downloads} 下）记 Seidel--Smith 引理 44、Abouzaid--Smith 2019 证明 $\mathrm{Kh}_{\rm symp}$ 与 Khovanov 同构——这是\textbf{纯数学}路线，不过 CS 路径积分。

\begin{keypoint}
\textbf{一句话}：1988 年的问题是“Jones 的三维定义是什么”；2000 年的问题是“怎样把 Jones 升级成同调”。后者更精细，但不回答前者。把两者放进同一章，只能作为“这篇论文开启的后续地图”，不能作为“这篇论文的内容”。
\end{keypoint}

\begin{boundary}
\begin{itemize}[leftmargin=1.4em]
\item 已确证：Euler 特征退回 Jones；Witten 2011 的 arXiv 号与 p.107 原句；Kapustin--Witten 的 arXiv 号；Abouzaid--Smith arXiv:1504.01230。
\item 未核：Khovanov 2000 的期刊卷页（Duke 101 (2000) 359--426 已核）；Khovanov 立方体构造的符号与 $q$ 次平移（未读原文）；“纯 CS 只有 ghost 数一维”是我们的重建。
\item 结论：本文不涉及范畴化；写成“延伸”时须整段标“后世”。
\end{itemize}
\end{boundary}
